%==============================================================================
\section{Thuật toán Mean-Shift}
\label{sec:mean-shift}
%==============================================================================

\begin{dinhnghia}[title={Mean-Shift Clustering}]
Mean-Shift là thuật toán phân cụm \textbf{không tham số hóa} (non-parametric): lặp dịch chuyển trung bình của điểm về vùng \textbf{mật độ cao nhất} của dữ liệu. Mật độ được định nghĩa qua \textbf{hàm nhân} (kernel), thường là Gaussian, gán trọng số theo khoảng cách tới mean hiện tại.
\end{dinhnghia}

\begin{ynghia}[title={Không cần chỉ định $K$ trước}]
Khác K-means phải cho biết số cụm, Mean-Shift \textbf{tự suy ra} số cụm từ dữ liệu và tham số bandwidth.
\end{ynghia}

\begin{tinhchat}[title={Phân cụm theo mật độ}]
Thuật toán thuộc nhóm density-based: cụm là vùng mật độ điểm cao, không chỉ dựa vào khoảng cách Euclidean đơn thuần giữa hai điểm xa nhau.
\end{tinhchat}

\subsection{Cách hoạt động Mean-Shift}

\begin{phuongphap}[title={Các bước}]
\begin{itemize}
  \item \textbf{Bước 1}: Ban đầu mỗi điểm (hoặc mỗi seed) được coi là một cụm riêng.
  \item \textbf{Bước 2}: Tính centroid (mean) có trọng số kernel quanh từng điểm.
  \item \textbf{Bước 3}: Cập nhật vị trí centroid mới.
  \item \textbf{Bước 4}: Lặp, dịch chuyển về vùng mật độ cao hơn.
  \item \textbf{Bước 5}: Dừng khi centroid không còn dịch chuyển đáng kể.
\end{itemize}
\end{phuongphap}

\subsection{Triển khai Mean-Shift trong Python}

\begin{phuongphap}[title={scikit-learn}]
Dùng \texttt{MeanShift} và \texttt{estimate\_bandwidth} trong \texttt{sklearn.cluster}.
\end{phuongphap}

\begin{vidu}[title={Import}]
\begin{lstlisting}
import numpy as np
import matplotlib.pyplot as plt
from sklearn.cluster import MeanShift, estimate_bandwidth
\end{lstlisting}
\end{vidu}

\begin{vidu}[title={Sinh dữ liệu và ước lượng bandwidth}]
\begin{lstlisting}
X = np.random.randn(500, 2)

bandwidth = estimate_bandwidth(X, quantile=0.1, n_samples=100)
\end{lstlisting}
\end{vidu}

\begin{ynghia}[title={Bandwidth}]
Bandwidth quyết định \textbf{độ rộng} kernel; ảnh hưởng trực tiếp số cụm và độ mịn của ranh giới.
\end{ynghia}

\begin{vidu}[title={Khởi tạo, huấn luyện, trực quan}]
\begin{lstlisting}
ms = MeanShift(bandwidth=bandwidth, bin_seeding=True)
ms.fit(X)

labels = ms.labels_
cluster_centers = ms.cluster_centers_
n_clusters_ = len(np.unique(labels))
print("Number of estimated clusters:", n_clusters_)

plt.figure(figsize=(7.5, 3.5))
plt.scatter(X[:, 0], X[:, 1], c=labels, cmap='viridis')
plt.scatter(cluster_centers[:, 0], cluster_centers[:, 1],
            marker='*', s=300, c='r')
plt.show()
\end{lstlisting}
\end{vidu}

\begin{vidu}[title={Ví dụ đầy đủ (dữ liệu Gaussian 2D)}]
\begin{lstlisting}
import numpy as np
import matplotlib.pyplot as plt
from sklearn.cluster import MeanShift, estimate_bandwidth

X = np.random.randn(500, 2)
bandwidth = estimate_bandwidth(X, quantile=0.1, n_samples=100)
ms = MeanShift(bandwidth=bandwidth, bin_seeding=True)
ms.fit(X)

labels = ms.labels_
cluster_centers = ms.cluster_centers_
print("Number of estimated clusters:", len(np.unique(labels)))

plt.scatter(X[:, 0], X[:, 1], c=labels, cmap='summer')
plt.scatter(cluster_centers[:, 0], cluster_centers[:, 1],
            marker='*', s=200, c='r')
plt.show()
\end{lstlisting}
\end{vidu}

\subsubsection{Ví dụ 2 --- Dataset 2D nhiều blob}

\begin{dongco}[title={make\_blobs 3D trong mã nguồn}]
Một số ví dụ dùng \texttt{make\_blobs} với tâm 3 chiều; có thể chiếu hai cột đầu để vẽ 2D hoặc dùng \texttt{centers=4} hai chiều tương đương bài minh họa 4 cụm.
\end{dongco}

\begin{vidu}[title={Sinh blob và vẽ}]
\begin{lstlisting}
from sklearn.datasets import make_blobs

X, _ = make_blobs(n_samples=700, centers=4,
                  cluster_std=0.5, random_state=0)
plt.scatter(X[:, 0], X[:, 1])
plt.show()
\end{lstlisting}
\end{vidu}

\begin{vidu}[title={Mean-Shift mặc định và in tâm}]
\begin{lstlisting}
ms = MeanShift()
ms.fit(X)
labels = ms.labels_
cluster_centers = ms.cluster_centers_
print(cluster_centers)
print("Estimated clusters:", len(np.unique(labels)))
\end{lstlisting}
\end{vidu}

\begin{vidu}[title={Vẽ từng điểm theo nhãn}]
\begin{lstlisting}
colors = 10 * ['r.', 'g.', 'b.', 'c.', 'k.', 'y.', 'm.']
for i in range(len(X)):
    plt.plot(X[i][0], X[i][1], colors[labels[i]], markersize=3)
plt.scatter(cluster_centers[:, 0], cluster_centers[:, 1],
            marker='.', color='k', s=20, linewidths=5, zorder=10)
plt.show()
\end{lstlisting}
\end{vidu}

\begin{tinhchat}[title={Kết quả mẫu}]
Ví dụ trên có thể ước lượng khoảng \textbf{3 cụm} với ba tâm cụm in ra từ \texttt{cluster\_centers\_}.
\end{tinhchat}

\subsection{Ứng dụng Mean-Shift}

\begin{ynghia}[title={Ứng dụng Mean-Shift}]
  \begin{itemize}
    \item Thị giác máy tính: Theo dõi đối tượng, phân đoạn ảnh, trích xuất đặc trưng.
    \item Xử lý ảnh: Phân đoạn ảnh theo độ tương đồng pixel.
    \item Phát hiện bất thường: Vùng mật độ thấp có thể đánh dấu bất thường.
    \item Phân khúc khách hàng và mạng xã hội: Gom nhóm hành vi hoặc sở thích người dùng tương tự.
  \end{itemize}
\end{ynghia}

\subsection{Ưu điểm và nhược điểm}

\begin{tinhchat}[title={Ưu điểm}]
\begin{itemize}
  \item Không giả định mô hình phân phối cố định như K-means hay GMM.
  \item Mô hình được cụm \textbf{lồi và phức tạp} (non-convex).
  \item Chỉ cần điều chỉnh chủ yếu tham số \textbf{bandwidth}; số cụm suy ra từ dữ liệu.
  \item Ít vấn đề \textbf{cục cực địa phương} như K-means.
  \item Tương đối ít bị ngoại lai kéo mean toàn cục như trung bình K-means.
\end{itemize}
\end{tinhchat}

\begin{luuy}[title={Nhược điểm}]
\begin{itemize}
  \item Chi phí tính toán cao trên dữ liệu \textbf{nhiều chiều}; số cụm có thể thay đổi đột ngột khi bandwidth không phù hợp.
  \item Không kiểm soát trực tiếp số cụm --- một số bài toán cần $K$ cố định.
  \item Khó phân biệt mode ``có ý nghĩa'' và mode nhiễu trên bề mặt mật độ.
\end{itemize}
\end{luuy}
